% !PRINTED 36
\weberpage{36}
%<<CONT
und da $U$ durch $\xi$, $V$ durch $\eta$ teilbar ist, so muß diese
Funktion durch $\xi\eta$ teilbar sein; setzen wir sie
$=\xi\eta(m\xi+n\eta)$, so ergibt sich aus (2)

\[
(a\xi+b\eta)^2=(\alpha^2\xi-\beta^2\eta)^2+\eta(m\xi+n\eta),
\]
\[
(a'\xi+b'\eta)^2=(\alpha^2\xi-\beta^2\eta)^2+\xi(m\xi+n\eta).
\]

Die Vergleichung der Koeffizienten ergibt

\[
a=\alpha^2, \qquad b=-\beta^2+\frac{m}{2\alpha^2}, \qquad b^2=\beta^4+n,
\]

\[
b'=\beta^2, \qquad a'=-\alpha^2+\frac{n}{2\beta^2}, \qquad a'^2=\alpha^4+m,
\]

und aus den beiden letzten Gleichungen jeder Reihe:

\[
m^2=4\alpha^2(n\alpha^2+m\beta^2),
\]
\[
n^2=4\beta^2(n\alpha^2+m\beta^2).
\]

Wenn man also

\[
m=h\alpha, \qquad n=h\beta
\]

setzt, so wird

\[
h=4\alpha\beta(\alpha+\beta).
\]

Hiernach lassen sich $a,b,a',b'$ durch die beiden Größen $\alpha,\beta$
ausdrücken in der Weise:

\weq{8}{\begin{aligned}
a&=\alpha^2, & b&=\beta^2+2\alpha\beta,\\
a'&=\alpha^2+2\alpha\beta, & b'&=\beta^2,
\end{aligned}}

und danach aus (4), (5)

\weq{9}{\varkappa^2=\frac{\beta^3}{\alpha^3}\,\frac{\beta+2\alpha}{\alpha+2\beta}, \qquad
\frac{\lambda}{\varkappa}=\frac{\alpha}{\beta}\,\frac{\beta+2\alpha}{\alpha+2\beta},}

oder durch Multiplikation und Division dieser beiden Gleichungen:

\weq{10}{\sqrt{\lambda\varkappa}=\frac{\beta}{\alpha}\,\frac{\beta+2\alpha}{\alpha+2\beta}, \qquad
\frac{\varkappa^3}{\lambda}=\frac{\beta^4}{\alpha^4}.}

Durch Elimination von $\beta:\alpha$ erhält man eine Gleichung zwischen
$\varkappa,\lambda$, die man die \gesperrt{Modulargleichung} nennt, deren Grad die
Anzahl der verschiedenen Transformationen dritten Grades angibt. Sie
nimmt die einfachste Gestalt an, wenn man setzt:

\weq{11}{\sqrt[4]{\varkappa}=u, \qquad \sqrt[4]{\lambda}=v.}

Dann werden die Gleichungen (10)

\[
u^2v^2=\frac{\beta}{\alpha}\,\frac{\beta+2\alpha}{\alpha+2\beta}, \qquad
\frac{u^3}{v}=\frac{\beta}{\alpha},
\]

also durch Einsetzen des Wertes von $\beta:\alpha$ in die erste
Gleichung und Beseitigung des Faktors $u^2$:

\weq{12}{v^4-u^4+2u^3v^3-2uv=0,}

eine Gleichung vom vierten Grade.
